% Page 007

\seite{7}

Beweise, so kann doch aus den oben geführten Thatsachen\ob
mit einiger Sicherheit auf eine vorübergegangene An\ohb
siedlung geschlossen werden, bis sie von nachrückenden\ob
germanischen Stämmen aus ihren Niederlassungen\ob
vertrieben wurden.

Eine andere Niederlassung gründeten die Römer in\ob
Sefferweich. Damit ist aber keineswegs behauptet, daß die\ob
\margmark{Funde aus\\der Zeit der\\Römer.\\um 1840 -\\1850}%
vorhandenen Niederlassungen vertrieben worden waren.\ob
Diese Niederlassung kann man durch folgende Thatsachen\ob
begründen: Vor einigen Jahren stieß ein Bauer beim\ob
Pflügen seines Feldes in \enquote{Bormeyer} rechts vom\ob
Weiher am Wege nach Seffern auf Widerstand.\ob
Er glaubte es mit einem Stein thun zu haben, als\ob
er denselben beseitigen wollte, stieß er auf ein\ob
Mauerwerk. In der frohen Hoffnung, einen großen\ob
Schatz zu finden, grub er nach und fand das Fundament\ob
eines         römischen Gebäudes. In demselben fand\ob
man Jagdtrophäen, eine Jagdpfeife, ein Hirschgeweih, eine\ob
eiserne Ginsterhacke von der Form eines Hufes eines\ob
Esels; ferner Münzen der Galba u Severus, Hirsch\ohb
fänger, Spieß, ein römisches Gewicht, bestehend aus\ob
einem pyramidisch behauenen Stein mit einem\ob
kupfernen Ringe ca 2 1/2 kg schwer. Die Kaminen, be\ohb
stehend aus großen Ziegeln, welche noch unversehrt\ob
unterirdisch durchgingen, weisen gleichzeitig auf ein\ob
römisches Bad hin. Dieses ist um so mehr anzunehmen,
\pend
